\documentclass[a4paper,12pt,oneside]{book} % Formato documento: foglio a4, font size 12, solo fronte
\usepackage[utf8]{inputenc} % Codifica caratteri
\usepackage[lmargin=2cm,rmargin=2cm,tmargin=2.5cm,bmargin=2cm]{geometry} % Bordi pagina come da formato tesi UNIBS
\linespread{1.5} % Interlinea
\setlength{\parindent}{0pt} % No spazio prima dei paragrafi
\setlength{\parskip}{1em}
\setlength{\headheight}{15pt}
\addtolength{\topmargin}{-2.5pt}

\usepackage{longtable}

\usepackage{graphicx} % Per poter inserire immagini
\usepackage{float} % Per controllare il posizionamento di immagini e tabelle
\usepackage{amsmath} % Per inserire equazioni
\usepackage{mathtools} % Per la generazione della graffa destra
\usepackage[hidelinks]{hyperref} % Per inserire URL
\usepackage{xurl} % Per consentire agli URL di spezzarsi a capo
\usepackage{enumitem} % Per poter inserire liste
\usepackage{siunitx} % Necessario per usare \setitemize
\setitemize{noitemsep,topsep=0pt,parsep=0pt,partopsep=0pt} % Per rendere le liste più compatte
\setenumerate{noitemsep,topsep=0pt,parsep=0pt,partopsep=0pt} % Per rendere le liste più compatte
\usepackage{booktabs} % Migliora qualità delle tabelle
\usepackage{multirow}
\usepackage{svg} % Inserimento di immagini SVG

\usepackage[titles]{tocloft}
\setlength{\cftbeforechapskip}{6pt}
\usepackage[backend=biber,sorting=none]{biblatex} % Per poter scrivere la bibliografia su un file a parte
\addbibresource{refs.bib} % File contenente la bibliografia

\usepackage{titlesec} % Per modificare il formato dei capitoli
\titleformat{\chapter}[hang]{\normalfont\Huge\bfseries}{\thechapter.}{8pt}{\Huge\bfseries}
\titlespacing*{\chapter}{0pt}{-40pt}{5pt}
\titlespacing*{\section}{0pt}{0pt}{0pt}
\titlespacing*{\subsection}{0pt}{0pt}{0pt}


\usepackage{fancyhdr} % Per modificare l'intestazione delle pagine
\pagestyle{fancy} % Customizzazione testata e piè di pagina
\renewcommand{\chaptermark}[1]{\markboth{#1}{#1}}
\fancyhf{}
\fancyhead[L]{\thechapter. \leftmark}
\fancyfoot[C]{\thepage}


\title{Progettazione sito Arcangelo DJ}
\author{Matteo Beatrice}
\date{Settembre 2022}

\begin{document}
    \begin{titlepage}
    \begin{figure}[H] % Logo Unibs
        \centering
        \includegraphics[width=72.4mm]{images/logo_unibs}
    \end{figure}

    \begin{center}
        \LARGE{\uppercase{Università degli Studi di Brescia}}\\ % Doppio backslash = vai a capo
        \vspace{5mm} % Spazio verticale
        \large{\uppercase{Dipartimento di Ingegneria dell'informazione}}\\
        \vspace{5mm}
        \large{Corso di Laurea Magistrale in Ingegneria Informatica}\\
    \end{center}

    \vspace{10mm}

    \begin{center}
        \LARGE{\textbf{Sistemi Informativi Evoluti e Big Data}}\\
    \end{center}

    \vspace{50mm}

    \begin{flushright}
        \large
        \textbf{Studente:}\\
        Matteo Beatrice (739848)
    \end{flushright}

    \vspace*{\fill} % Spazio bianco per mettere il testo in fondo alla pagina

    \rule{0.8\textwidth}{0.6pt}\\ % Linea nera
    \centering{\Large{Anno Accademico 2022/2023}}
\end{titlepage} %Frontespizio
    \hyphenpenalty=2000 % Rende meno probabile lo spezzamento delle parole a capo

    \frontmatter
    \pagenumbering{Roman} % Numeri romani maiuscoli per numero pagina per indice e intro
    \tableofcontents % Indice
    \chapter{Introduzione}

In questo report sono presenti tutte le specifiche di progettazione necessarie per la realizzazione del sito Arcangelodj.it\\
Il capitolo~\ref{ch:requirements} contiene la Specifica dei Requisiti con la descrizione degli utenti che interagiranno con il sito e la descrizione della gerarchia e delle funzionalità.\\
 % Introduzione

    \mainmatter
    \chapter{Specifica dei Requisiti}\label{ch:requirements}
In questo capitolo vengono descritti:
\begin{itemize}
  \item Gerarchia degli utenti che interagiranno con il sito.
  \item Requisit funzionali degli utenti
  \item Dizionario dei Dati
  \item Site View
\end{itemize}

\section{Gerarchia degli Utenti}\label{sec:user_hierarchy}

\begin{figure}[H]
  \centering
  \includesvg[inkscapelatex=false, scale=0.60]{diagrams/User Hierarchy.svg}
  \caption{User Hierarchy}
\end{figure}

\section{Schede Utenti}\label{sec:users}

Dei 3 casi d'uso specificati viene presentata solamente la scheda di specifica dell'accesso alla lista 
degli eventi futuri perché gli altri due riguardano solamente la visualizzazione di una pagina statica:
\begin{itemize}
    \item \textbf{Accedi a Galleria} $\rightarrow$ Utente visualizza la pagina della galleria con foto e video di eventi passati
    \item \textbf{Accedi a Contatti e Membri del Team} $\rightarrow$ Utente visualizza i diversi canali attraverso i quali può contattare il team ed una lista di tutti i membri con foto e ruolo all'interno del team
\end{itemize}

\begin{table}[H]
    \begin{center}  
        \begin{tabular}{ | l | p{13cm} |} % you can change the dimension according to the spacing requirements  
        \hline
        \textbf{Titolo} & \textbf{Accedi a lista Eventi Futuri}\\ \hline  
           
        Scopo & Come gli Utenti Esterni non Registrati possono visualizzare e filtrare una lista ordinata dei prossimi eventi programmati ed i relativi dettagli\\ \hline  
           
        Pre-Condizione & Nessuna\\ \hline  
           
        Post-Condizione & Utente visualizza correttamente gli eventi (ordinati in base alla data) ed i dettagli eventualmente filtrando la lista in base alla venue desiderata o in base a del testo inserito nella search box\\ \hline  
           
        Workflow & Il caso d'uso prevede i seguenti passi
        \begin{itemize}
            \item Visualizza una versione limitata della lista (i prossimi 4/5 eventi) direttamente nella homepage oppure entra nella pagina dedicata
            \item Nella pagina dedicata l'utente visualizza la lista di tutti gli eventi ed imposta i filtri
            \item Lista viene aggiornata e utente visualizza eventi che rispettano i filtri imposti
        \end{itemize}\\  
        \hline
           
        \end{tabular}  
        \caption{Scheda di Specifica - Accesso alla lista degli eventi}
    \end{center}  
\end{table}
Non viene presentata la scheda di specifica del Log In e della Modifica dei Dati del profilo perché sono casi d'uso comuni

\begin{itemize}
    \item \textbf{Log In} $\rightarrow$ Utente effettua autenticazione utilizzando username(email) e password forniti quando ha effettuato la registrazione
    \item \textbf{Modifica dei dati del profilo} $\rightarrow$ Utente ha la possibilità di accedere alla sua pagina personale e modificare i propri dati
\end{itemize}

\begin{table}[H]
    \begin{center}  
        \begin{tabular}{ | l | p{13cm} |} % you can change the dimension according to the spacing requirements  
        \hline
        \textbf{Titolo} & \textbf{Crea Prenotazione}\\ \hline  
           
        Scopo & Quello che deve fare l'utente per poter creare una nuova prenotazione ad un evento futuro al quale è interessato\\ \hline  
           
        Pre-Condizione & Utente effettua il Log In con la propria utenza e visualizza i dettagli dell'evento al quale è interessato\\ \hline  
           
        Post-Condizione & La nuova prenotazione per l'evento scelto viene creata rispettando i parametri definiti dall'utente\\ \hline  
           
        Workflow & Il caso d'uso prevede i seguenti passi
        \begin{itemize}
            \item Utente effettua il Log In
            \item Trova un evento al quale è interessato e visualizza i dettagli
            \item Essendo un utente registrato ha la possibilità di creare una prenotazione
            \item Utente visualizza un form dove gli vengono chiesti i dati relativi alla prenotazione
            \item Utente conferma i dati inseriti e se i posti ancora disponibili per l'evento sono sufficienti la prenotazione viene creata e confermata
        \end{itemize}\\ \hline  
           
        \end{tabular}  
        \caption{Scheda di Specifica - Creazione di una prenotazione}
    \end{center}  
\end{table}

\begin{table}[H]
    \begin{center}  
        \begin{tabular}{ | l | p{13cm} |} % you can change the dimension according to the spacing requirements  
        \hline
        \textbf{Titolo} & \textbf{Modifica Prenotazione}\\ \hline  
           
        Scopo & Quello che deve fare l'utente per poter effettuare delle modifiche alla prenotazione\\ \hline  
           
        Pre-Condizione & Utente ha effettuato il Log In e la prenotazione da modificare deve essere per un evento futuro\\ \hline  
           
        Post-Condizione & Le eventuali modifiche vengono salvare o la prenotazione viene cancellata se l'utente ha deciso di annullarla\\ \hline  
           
        Workflow & Il caso d'uso prevede i seguenti passi
        \begin{itemize}
            \item Utente effettua il Log In
            \item Attraverso la lista delle sue prenotazioni visualizza i dettagli di una sua prenotazione
            \item Effettua le modifiche oppure sceglie di annullare la prenotazione
            \item Se modifiche rispettano il numero di posti ancora disponibili per l'evento la prenotazione viene aggiornata oppure viene annullata(nel caso di richiesta di annullamento)
            \item Utente visualizza le proprie modifiche nella pagina delle prenotazioni
        \end{itemize}\\ \hline  
           
        \end{tabular}  
        \caption{Scheda di Specifica - Modifica di una prenotazione}
    \end{center}  
\end{table}
\begin{table}[H]
    \begin{center}  
        \begin{tabular}{ | l | p{10cm} |} % you can change the dimension according to the spacing requirements  
        \hline
        \textbf{Nome del Gruppo} & \textbf{Gestore Venue}\\ \hline  
           
        Descrizione & Utenti che gestiscono le Venue dove si terranno gli Eventi e sono interessati alla sottomissione di richieste di creazione/modifica di nuovi Eventi e alla visualizzazione delle Prenotazioni riguardanti Eventi che si tengono nelle loro Venue\\ \hline  
           
        Dati di Profilo & Nome e Cognome, username/email e password per il login.\\ \hline  
           
        Oggetti Accessibili in Lettura & Oggetti accessibli dalla parte pubblica, Venue gestite ed Eventi collegati con relative Prenotazioni presenti.\\ \hline  
           
        Oggetti Accessibili in Scrittura & Modifiche al proprio profilo, creazione/modifica di Venue e sottomissione di richieste di creazione/modifica di Eventi.\\ \hline  
           
        \end{tabular}  
        \caption{Scheda Gestore di Venue}
    \end{center}  
\end{table}
\begin{table}[H]
    \begin{center}  
        \begin{tabular}{ | l | p{10cm} |} % you can change the dimension according to the spacing requirements  
        \hline
        \textbf{Nome del Gruppo} & \textbf{Membro del Team}\\ \hline  
           
        Descrizione & Utente Interno che gestisce la galleria e la creazione di nuovi eventi (accettando le richieste dei \textit{Gestori delle Venue})\\ \hline  
           
        Dati di Profilo & Nome e Cognome, username/email e password per il login. Una fotografia e un ruolo che saranno visualizzati anche nella parte pubblica del sito (Sezione membri del team).\\ \hline  
           
        Oggetti Accessibili in Lettura & Oggetti accessibli dalla parte pubblica, Venue, Eventi (e richieste di modifica/creazione) e Prenotazioni.\\ \hline  
           
        Oggetti Accessibili in Scrittura & Modifiche al proprio profilo, aggiunta/modifica di oggetti nella galleria, creazione/modifica di Venue, Eventi(e Richieste) e Prenotaizoni.\\ \hline  
           
        \end{tabular}  
        \caption{Scheda Membro del Team}
    \end{center}  
\end{table}
\begin{table}[H]
    \begin{center}  
        \begin{tabular}{ | l | p{10cm} |} % you can change the dimension according to the spacing requirements  
        \hline
        \textbf{Nome del Gruppo} & \textbf{Amministratore}\\ \hline  
           
        Descrizione & Membro del Team che ha anche la possibilità di "promuovere" profili utenti base a \textit{Membro del Team} e \textit{Gestore Venue}\\ \hline  
           
        Dati di Profilo & Nome e Cognome, username/email e password per il login.\\ \hline  
           
        Oggetti Accessibili in Lettura & Uguale al \textit{Team Member} ma visualizza anche una lista degli utenti registrati.\\ \hline  
           
        Oggetti Accessibili in Scrittura & Uguale al \textit{Team Member} ma può modificare l'elevazione degli utenti registrati.\\ \hline  
           
        \end{tabular}  
        \caption{Scheda Amministratore}
    \end{center}  
\end{table}

\newpage
\section{Casi d'Uso}\label{sec:use_cases}
\subsection{Utente Esterno non Registrato}\label{subsec:uc_unregistered}

\begin{figure}[H]
    \centering
    \includesvg[inkscapelatex=false, scale=0.60]{diagrams/UC Utente Esterno non Registrato.svg}
    \caption{Casi d'uso per l'utente esterno non registrato}
\end{figure}

Dei 3 casi d'uso specificati viene presentata solamente la scheda di specifica dell'accesso alla lista 
degli eventi futuri perché gli altri due riguardano solamente la visualizzazione di una pagina statica:
\begin{itemize}
    \item \textbf{Accedi a Galleria} $\rightarrow$ Utente visualizza la pagina della galleria con foto e video di eventi passati
    \item \textbf{Accedi a Contatti e Membri del Team} $\rightarrow$ Utente visualizza i diversi canali attraverso i quali può contattare il team ed una lista di tutti i membri con foto e ruolo all'interno del team
\end{itemize}

\begin{table}[H]
    \begin{center}  
        \begin{tabular}{ | l | p{13cm} |} % you can change the dimension according to the spacing requirements  
        \hline
        \textbf{Titolo} & \textbf{Accedi a lista Eventi Futuri}\\ \hline  
           
        Scopo & Come gli Utenti Esterni non Registrati possono visualizzare e filtrare una lista ordinata dei prossimi eventi programmati ed i relativi dettagli\\ \hline  
           
        Pre-Condizione & Nessuna\\ \hline  
           
        Post-Condizione & Utente visualizza correttamente gli eventi (ordinati in base alla data) ed i dettagli eventualmente filtrando la lista in base alla venue desiderata o in base a del testo inserito nella search box\\ \hline  
           
        Workflow & Il caso d'uso prevede i seguenti passi
        \begin{itemize}
            \item Visualizza una versione limitata della lista (i prossimi 4/5 eventi) direttamente nella homepage oppure entra nella pagina dedicata
            \item Nella pagina dedicata l'utente visualizza la lista di tutti gli eventi ed imposta i filtri
            \item Lista viene aggiornata e utente visualizza eventi che rispettano i filtri imposti
        \end{itemize}\\  
        \hline
           
        \end{tabular}  
        \caption{Scheda di Specifica - Accesso alla lista degli eventi}
    \end{center}  
\end{table}
Non viene presentata la scheda di specifica del Log In e della Modifica dei Dati del profilo perché sono casi d'uso comuni

\begin{itemize}
    \item \textbf{Log In} $\rightarrow$ Utente effettua autenticazione utilizzando username(email) e password forniti quando ha effettuato la registrazione
    \item \textbf{Modifica dei dati del profilo} $\rightarrow$ Utente ha la possibilità di accedere alla sua pagina personale e modificare i propri dati
\end{itemize}

\begin{table}[H]
    \begin{center}  
        \begin{tabular}{ | l | p{13cm} |} % you can change the dimension according to the spacing requirements  
        \hline
        \textbf{Titolo} & \textbf{Crea Prenotazione}\\ \hline  
           
        Scopo & Quello che deve fare l'utente per poter creare una nuova prenotazione ad un evento futuro al quale è interessato\\ \hline  
           
        Pre-Condizione & Utente effettua il Log In con la propria utenza e visualizza i dettagli dell'evento al quale è interessato\\ \hline  
           
        Post-Condizione & La nuova prenotazione per l'evento scelto viene creata rispettando i parametri definiti dall'utente\\ \hline  
           
        Workflow & Il caso d'uso prevede i seguenti passi
        \begin{itemize}
            \item Utente effettua il Log In
            \item Trova un evento al quale è interessato e visualizza i dettagli
            \item Essendo un utente registrato ha la possibilità di creare una prenotazione
            \item Utente visualizza un form dove gli vengono chiesti i dati relativi alla prenotazione
            \item Utente conferma i dati inseriti e se i posti ancora disponibili per l'evento sono sufficienti la prenotazione viene creata e confermata
        \end{itemize}\\ \hline  
           
        \end{tabular}  
        \caption{Scheda di Specifica - Creazione di una prenotazione}
    \end{center}  
\end{table}

\begin{table}[H]
    \begin{center}  
        \begin{tabular}{ | l | p{13cm} |} % you can change the dimension according to the spacing requirements  
        \hline
        \textbf{Titolo} & \textbf{Modifica Prenotazione}\\ \hline  
           
        Scopo & Quello che deve fare l'utente per poter effettuare delle modifiche alla prenotazione\\ \hline  
           
        Pre-Condizione & Utente ha effettuato il Log In e la prenotazione da modificare deve essere per un evento futuro\\ \hline  
           
        Post-Condizione & Le eventuali modifiche vengono salvare o la prenotazione viene cancellata se l'utente ha deciso di annullarla\\ \hline  
           
        Workflow & Il caso d'uso prevede i seguenti passi
        \begin{itemize}
            \item Utente effettua il Log In
            \item Attraverso la lista delle sue prenotazioni visualizza i dettagli di una sua prenotazione
            \item Effettua le modifiche oppure sceglie di annullare la prenotazione
            \item Se modifiche rispettano il numero di posti ancora disponibili per l'evento la prenotazione viene aggiornata oppure viene annullata(nel caso di richiesta di annullamento)
            \item Utente visualizza le proprie modifiche nella pagina delle prenotazioni
        \end{itemize}\\ \hline  
           
        \end{tabular}  
        \caption{Scheda di Specifica - Modifica di una prenotazione}
    \end{center}  
\end{table}

    % \chapter{Kernel Search}\label{ch:kernel-search}
Kernel Search is a heuristic framework used to solve MIP problems, which was introduced in~\cite{kernel:2007}.\\
The algorithm has been studied in a few applications, including the \textit{Multidimensional Knapsack Problem}
in~\cite{kernel:2010} and the \textit{Portfolio Selection Problem} in~\cite{kernel:2012}.

The algorithm is based on a few observations that can often be made when solving MIP problems:
\begin{itemize}
    \item In the optimal solution there are only a few non-zero variables;
    \item Basic variables in the LP-relaxation are good predictors of non-zero variables in the MIP optimum;
    \item Reduced costs are good predictors of the likelihood of a non-basic variable to be in the MIP optimum.
\end{itemize}

Kernel search needs an MIP solver such as GUROBI or CPLEX to solve a set of MIP sub-problems.


\section{Algorithm}
The algorithm has two main phases: an \textit{initialization phase}
that builds an initial \textit{kernel set} and a number of \textit{buckets},
and an \textit{improvement phase} that iteratively enlarges the kernel set and improves the solution.

\subsection{Initialization Phase}
The first step of the Kernel Search is to solve the LP relaxation of the problem.\\
Then, the set of variables of the model is sorted according to some \textit{sorting criterion}.

Once this is done, the kernel set \(\Lambda\) is built by selecting the \textit{C}
variables in the sorted set according a \textit{kernel construction criterion},
where \textit{C} is the size of the kernel set.

The next step is to solve MIP(\(\Lambda\)),
which is the original problem restricted to only include the variables in \(\Lambda\):
all the variables not in \(\Lambda\) are set to 0 to exclude them from the problem.\\
The resulting solution will be called \(x^{*}\) and the optimum value \(w^{*}\).\\
It is possible that no solution is found for MIP(\(\Lambda)\): in this case \(w^{*}\) is set to \(-\infty\).

The variables not in \(\Lambda\) are partitioned into a certain number of buckets \textit{nb}, according
to a \textit{bucket construction criterion}.

\subsection{Improvement Phase}
The main objectives of this phase are finding a new improving feasible solution and identifying new variables
to enter the kernel set.

Once the buckets are constructed, the algorithms proceeds with the improvement phase:\\
for each bucket \(B_{i} \ (i=1,\dots,nb)\),
the restricted sub-problem MIP(\(\Lambda \cup B_{i}\)) is solved.

Only the feasible sub-problems with an incumbent solution \(\bar{x^{*}} \geq w^{*}\) are considered.\\
For such sub-problems, the set \(\bar{\Lambda}_{i} \coloneqq \{\text{selected variables in } B_{i}\}\) is built.\\
Then the kernel set is updated: \(\Lambda \coloneqq \Lambda \cup \bar{\Lambda}_{i}\).
Since variable are only added to the kernel, and never removed, its size increases monotonically.

After all the \textit{nb} have been iterated, the algorithm ends.


\section{Parameters}
The Kernel Search has quite a few configuration parameters that can be used to tune the algorithm for a specific
problem:
\begin{itemize}
    \item \textit{Variable sorting criterion}: used to sort the variables in the initialization phase.
    Ideally, if the sorting criterion is selected appropriately,
    all the significant variables for the problem should be located in \(\Lambda\) and the very first buckets.
    \item \textit{Kernel size C}: the number of items initially selected to construct the kernel.
    \item \textit{Kernel construction criterion}: used to build the kernel.
    The criterion defined in the basic Kernel Search consists in selecting
    the first C variables in the sorted set, but for certain problems more complex criteria can be used.
    \item \textit{Bucket construction criterion}: how the buckets are constructed.
    Also defines the number and size of the buckets.
\end{itemize}


\section{Improving Efficiency of the Basic Kernel Search}\label{sec:improving-efficiency}
To improve the performance of the basic Kernel Search, two constraints can be added
when solving each MIP(\(\Lambda \cup B_{i}\)):
\begin{equation}
    \label{eq:impr1}
    \text{value of the objective function} \geq w^{*}
\end{equation}
\begin{equation}
    \label{eq:impr2}
    \sum_{j \in B_{i}} x_{j} \geq 1
\end{equation}
The \textit{cutoff constraint}~\eqref{eq:impr1} allows only solutions that improve on the current
objective value.
Constraints~\eqref{eq:impr2} ensure that at least one item must be selected from the bucket \(B_{i}\):
such an item will then be included in the new kernel set.


\section{Iterative Kernel Search}\label{sec:iter}
Iterative Kernel Search is a variation of the basic Kernel Search where buckets are scrolled more than once.
The algorithm is as follows:
\begin{enumerate}
    \item Execute the basic Kernel Search;
    \item Set \(nb \coloneqq q-1\), where \(q \coloneqq \max\{i:\bar{\Lambda}_{i} \neq \emptyset\}\);
    \item Execute the improvement phase of the algorithm;
    \item If a new variable enters the kernel set, \(nb\) is reset to its initial value,
    the improvement phase is executed again, and the algorithm is repeated from step 2.
\end{enumerate}

A simpler version of this algorithm repeats step 3 a certain number of times,
skipping steps 2 and 4 entirely.
The \textit{number of iterations} is a configuration parameter.
    % \chapter{Implementation}\label{ch:implementation}
The Java source code for the Kernel Search was provided during the course,
and it implements a simple iterative Kernel Search (as explained in~\ref{sec:iter}),
using GUROBI as the MIP solver.
It also includes the improvements described in~\ref{sec:improving-efficiency}.

The code, which can be found alongside its documentation at
\url{https://github.com/Golino98/KernelSearchGolinoCottiBeatrice},
was initially modified to implement the MKP model described in~\ref{sec:model},
and then refactored to improve the code quality and efficiency.\\
The algorithm itself and the configuration parameters were not changed at this time.

The Kernel Search is configured with the following parameters:
\begin{itemize}
    \item \textit{Variable sorting criterion}: sort variables by non-increasing value and non-decreasing reduced cost.
    \item \textit{Kernel size C}: 15\% of the number of variables (rounded to the nearest integer).
    \item \textit{Bucket construction criterion}: iterates through the sorted variables, grouping them in
    buckets of size equal to the 2.5\% of the number of variables (rounded to the nearest integer).
    It's possible for the last bucket to contain fewer items than the previous ones, because the number
    of remaining variables could be inferior to the bucket size.
    \item \textit{Number of iterations}: 2.
\end{itemize}


\section{GUROBI}
The project uses GUROBI 9.5, which is the latest version of GUROBI available at the time of writing.\\
The GUROBI configuration is the following:
\begin{itemize}
    \item \textit{Presolve}: 2 (aggressive presolve)
    \item \textit{MIPGap}: \num{1e-12}
    \item \textit{Threads}: 12
\end{itemize}

    % \chapter{Computational Experiments}\label{ch:computational-experiments}
In this section we report the results of tests executed on the default Kernel Search applied to the MKP\@.
The main objective of these experiments is to determine the performance of the default Kernel Search,
both in terms of quality of the solution found and time required for the execution.


\section{Test Instances}
The test instances for the MKP were kindly provided to us by the authors of~\cite{mkp:2019}.\\
There are in total 2100 test instances, organized in
five directories of increasing complexity:
\begin{itemize}
    \item SMALL
    \item FK\_1
    \item FK\_2
    \item FK\_3
    \item FK\_4
\end{itemize}
SMALL contains 180 instances with \(m \in {10,20} \) and
\(n \in {20,40,60}\), while FK\_1, FK\_2, FK\_3, FK\_4 contain 480
instances each, designed with the aim of identifying critical ratios of
\(n/m\) that produce difficult instances.

The weights \(w_{j}\) are always uniformly distributed in \([\alpha, 1000]\), with \(\alpha = 1\) for the
SMALL instances and \(\alpha = 10\) for the FK instances.

Each instance belongs to one of the following four correlation classes:
\begin{itemize}
    \item \textit{uncorrelated}: profits \(p_{j}\) are uniformly distributed in \([\alpha,1000]\);
    \item \textit{weakly correlated}: For SMALL, \(p_{j} = 0.6 w_{j} + \theta_{j}\) , with \(\theta_{j}\) uniformly random in
    \([1, 400]\).
    For FK, the \(p_{j}\) values are uniformly distributed in \([w_{j} - 100, w_{j} + 100]\), such that \(p_{j} \geq 1\);
    \item \textit{strongly correlated}: For SMALL, \(p_{j} = w_{j} + 200\), for FK, \(p_{j} = w_{j} + 10\);
    \item \textit{subset-sum}: \(p_{j} = w_{j}\).
\end{itemize}

All the instances in SMALL and FK\_1 were already solved to the optimum.
Of those in the three remaining folders only a few were already solved:
the unsolved ones only had an upper and lower bound.

\subsection{Format of the Instances}
The instances are contained in plain \texttt{.txt} files.

An example of the structure of an instance is shown below:
\[
    \begin{split}
        &\begin{rcases*}
             3 \quad
        \end{rcases*}
        \rightarrow \text{integer that indicates the value of \textit{m (number of knapsacks)}.}\\
        &\begin{rcases*}
             7 \quad
        \end{rcases*}
        \rightarrow \text{integer that indicates the value of \textit{n (number of items)}.}\\
        &\begin{rcases*}
             908 \quad\\
             834\\
             675\\
        \end{rcases*}
        \rightarrow \text{integers that indicate the capacity } c_{i} \text{ of the knapsack i with }
        \space i=1,\dots,m.\\
        &\begin{rcases*}
             264    \quad 430 \quad\\
             606 \quad    945\\
             268 \quad    409\\
             619 \quad    591\\
             958 \quad    839\\
             972 \quad    818\\
             723 \quad    71\\
        \end{rcases*}
        \rightarrow
        \begin{rcases*}
            \text{First column shows the weight } w_{i} \quad \\
            \text{Second column shows the profit } p_{i}.
        \end{rcases*}
        \rightarrow
        \space i=1,\dots,n.
    \end{split}
\]

The execution times required to solve the instances to the optimum were not provided to us.
For a few of the instances we were able to find the exact solution
using a GUROBI solver with the following configuration:

\begin{itemize}
    \item \textit{Presolve}: 2 (aggressive presolve)
    \item \textit{Time limit}: 1 hour
    \item \textit{Threads}: 12
    \item \textit{MIPGap}: \num{1e-12}
\end{itemize}

Most of the instances however used all the available time without reaching the optimum.
The code and the results can be found at \url{https://github.com/Golino98/EsattoMKP}.

\subsection{Selected Instances}\label{subsec:inst}
To simplify the testing of the changes to the Kernel Search algorithm,
we only considered 40 instances: 6 solved to the optimum for each directory,
plus 10 that were not solved to the optimum.

Tables~\ref{tab:exact_opt} and~\ref{tab:exact_noopt} represent,
respectively, the optimum values of the 30 instances solved to
the optimum, and the upper and lower bound for the 10
that couldn't be solved.

\input{tables/exact_opt.tex}
\input{tables/exact_noopt.tex}


\section{Performance of the default Kernel Search}
Table ~\ref{tab:ker_per_val_abs_RC} outlines the result of tests executed on
the Kernel Search algorithm with the default configuration,
as outlined in ~\ref{ch:implementation}.
The tables include, for each instance:
\begin{itemize}
    \item Directory of the instance;
    \item Name of the instance;
    \item Solution found;
    \item Time elapsed;
    \item A boolean value that specifies if the time limit (120s) was reached.
\end{itemize}
\input{tables/kernelBase/KernelBuilderPercentage/VariableSorterByValueAndAbsoluteRC}
    % \chapter{Improvements}\label{ch:improvements}
In this section we will describe in chronological order the various changes
that we applied to the default Kernel Search implementation,
with the objective of improving its performance
in solving the MKP instances.

For the sake of completeness, we will describe all attempts at improving
the algorithm, including the unsuccessful ones.


\section{Variable Sorting}
As outlined in~\cite{kernel:2010} and~\cite{kernel:2012} the quality of
the solutions found by Kernel Search heavily depends on the criterion used to sort the variables.

In ~\cite{kernel:2010}, a sorting criterion for the \textit{Multidimensional Knapsack Problem} is mentioned:
sorting by non-increasing \textit{LP relaxation value} and non-decreasing \textit{absolute value of the reduced cost}.
This is the default criterion that was provided in the code, and
the experiments we executed on the instances proved that this sorting also
provides high quality solutions for the MKP\@.

Other criteria we have tested are:
\begin{itemize}
    \item Sort by non-decreasing \textit{reduced cost}, breaking ties using the non-increasing {LP relaxation value};
    \item Sort by non-increasing \textit{ratio of profit and weight of the item}, breaking ties using the
    non-increasing \textit{reduced cost value};
    \item Sort by non-increasing \textit{value of the LP relaxation} multiplied by the \textit{ratio of profit and
    weight},    breaking ties using the non-decreasing \textit{reduced cost};
    \item Sort randomly.
\end{itemize}

Generally, the efficiency of a sorting criterion depends on the kernel construction criterion used.


\section{Kernel Construction Criterion}\label{sec:kernel_criteria}
A recurring problem we have identified when running complex
instances (such as the one in FK\_4) is that,
from the very beginning, the kernel contains quite a lot of variables,
which makes the solving of the kernel problem and the buckets sub-problems slow.

This means that a starting point to improve the performance of the Kernel Search
could be to optimize the management of the kernel set,
possibly starting from the \textit{kernel construction criterion},
which by default is to include the first \textit{C} sorted variables
(Percentage kernel).

We have experimented with the following criteria:
\begin{enumerate}
    \item Positive kernel: select the variables that have a value greater than 0 in the LP relaxation;
    \item Integer kernel: select the variables that have a value equal to 1 in the LP relaxation;
    \item Threshold kernel: select the variables that have a LP relaxation value greater than a threshold,
    that we have set to 0.6.
\end{enumerate}


\section{Overlapping Buckets}
A criterion that could improve the quality of the solution of the bucket
sub-problems is to use partially overlapping buckets:
this could potentially increase the possibility of related
variables to be included together in the kernel set.

We implemented this feature as an alternative bucket construction criterion.
This feature works well with certain combinations of
kernel construction criterion and variable sorting criterion.


\section{Ejection of Variables from the Kernel}
Another possible improvement for the efficiency of the Kernel Search,
is the \textit{kernel eject} method, which allows to remove variables from the kernel.
This technique is particularly useful for more complex instances, where every
bucket sub-problem normally takes a significant amount of time to be solved.

To implement this functionality, at each iteration of the Kernel Search
we keep a counter \(h\) for the number of solutions found.
Each variable that isn't in the initial kernel
is also associated with a counter \(k_{v}, v \in variables \setminus\) \textit{variables in the initial kernel}
that is incremented when the variable has a non-zero value in the solution of a bucket sub-problem.
At the end of each iteration, for each variable \(v\), if
\((h - k_{v}) - k_{v} <= threshold\), the variable is removed from the kernel.

The values of \(h\) and \(k_{v}\) are reset at each iteration of the Kernel Search.

The \textit{eject threshold} is a configuration parameter:
a low threshold will remove more variables from the kernel, increasing efficiency at the cost of the quality of the
solution found.
A high threshold instead has a more subtle effect.

When the variables are removed from the kernel they are not deleted from the model,
but they are returned to their original bucket.

Experiments on the instances demonstrated that this method significantly cuts
the solving time of the bucket sub-problems.
The quality of the solution found is worsened, but on the flip side
the more iterations and buckets can be solved in the same time limit.


\section{Repetition Counter}
An observation we could make is that when solving buckets,
after a certain number of iterations the algorithm
repeatedly finds new solutions with the same objective value,
and sometimes it even struggles to find new solutions at all.
In other words, the algorithm gets stuck in \textit{local optima},
from which it's hard to escape.

Our hypothesis is that there are two causes for this:
\begin{enumerate}
    \item The GUROBI solver cannot find a (better) solution for the sub-problem,
    either due to the infeasibility of the problem or because of
    the time limit set for solving each bucket.
    \item The Kernel Search algorithm, during its improvement phase,
    only accepts a new solution (thus updating the kernel)
    if it improves upon, or is at least equal, to the incumbent one.
    This is done by adding a cutoff constraint, as explained in~\ref{sec:improving-efficiency}.
\end{enumerate}

In order to mitigate this problem we introduced a \textit{repetition counter}
that, during the improvement phase of Kernel Search,
removes the cutoff constraint for \textit{k} buckets
when the same solution (or no solution) is found for \textit{h} times.

The idea is that this method allows to select if the focus should be on
\textit{diversification} or \textit{intensification},
by appropriately setting parameters \textit{h} and \textit{k}.

A low value for \textit{h} and a high one for \textit{k} allow for
diversification, by adding to the kernel variables that otherwise
may never be selected, and could allow escaping the local optimum.

On the opposite, a high value for \textit{h} and a low one for \textit{k}
switch to focus on intensification, by giving priority to finding
better solutions.

After experimenting with different values for \textit{h} and \textit{k},
we found that keeping them more or less equal allowed for
a reasonable balance between diversification and intensification.
In particular we found that \(h=3\) and \(k=3\) worked quite well
with the test instances we selected.

\subsection{Reset counter on new optimum}
A small enhancement we have applied is to reset the counter to its initial status
whenever the Kernel Search finds a new solution
that is better than any other found before.

This is significant because the default repetition counter
method completely ignores the value of the solutions found
during the \textit{k} cycles.


\section{Item Dominance}
In~\cite{mkp:2019} an improvement for the MKP model is suggested:
given two items \textit{j} and \textit{k}, if \(w_{j} \leq w_{k}\)
and \(p_{j} \geq p_{k}\), then it is said that \textit{j dominates k}.
This means that when an item is excluded from the solution,
all items dominated by it can also be excluded.

The simplest way to implement this method is to preliminarily sort all
items according to non-increasing weight, breaking ties by
non-decreasing profit.
For each item \(k\), items \(j \coloneqq k+1,\dots,n\) are parsed,
and if \(p_{j} \geq p_{k}\) then the pair \((j,k)\) is added to a
dominance list \(D\).
Then, the following constraints are added to the model:
\begin{equation}
    \label{eq:itemdom}
    \sum_{i=1}^{m} x_{ij} \geq \sum_{i=1}^{m} x_{ik} \qquad (j,k) \in D
\end{equation}

To efficiently implement this technique,the dominance list is only built once at startup
and the constraints~\eqref{eq:itemdom} are applied when solving
the relaxation, the kernel problem and the bucket sub-problems.

The testing proved that this method increases performance of the Kernel Search with very little
overhead costs.


\section{Instance Reduction}
Another improvement for the MKP model describe in~\cite{mkp:2019} is
\textit{instance reduction}.

Let \textit{I} be any subset of knapsacks and let \textit{J} be the set of
all items of the instance that can be packed in a knapsack of \textit{I}:
\begin{equation}
    J \coloneqq \{j:w_{j} \leq \max_{i \in I} \{c_{i}\}, 1 \leq j \leq n\}
\end{equation}
If there exists a feasible packing of the items of \textit{J} into the knapsacks
of \textit{I}, then such packing can be fixed, and sets \textit{I} and \textit{J}
can be removed from the instance.

The most efficient way to implement this property is to start
by sorting the knapsacks by non-decreasing capacity, add the first knapsack
to set \textit{I}, and iteratively add the next (smaller) knapsack to \textit{I}.
At each iteration, we add the appropriate items to \textit{J},
and check if \(\sum_{j \in J} w_{j} \leq \sum_{i \in I} c_{i}\).
This allows to avoid running the more expensive bin packing algorithm
when it's guaranteed that there is no feasible solution.

To solve the bin packing problem, in~\cite{mkp:2019} an exact method is suggested.
However, the execution time for such an algorithm is unacceptable for this project.
We decided anyway to try solving the bin packing problem using a modified
version of the \textit{First-Fit-Decreasing (FFD)} heuristic.

The basic version of FFD works as follows:
Given \textit{n} items with a weight and a fixed capacity for each bin:
\begin{enumerate}
    \item Order the items from largest to smallest;
    \item Open a new empty bin;
    \item For each item, from largest to smallest, find the first bin into which the item fits, if any.
    If such a bin is found, put the item in it.
    Otherwise, open a new empty bin and put the item in it.
\end{enumerate}

In this algorithm the bins correspond to the knapsacks of \textit{I},
and the items are the ones contained in \textit{J}.
Also, the number of bins is pre-determined by the size of \textit{I},
and each bin has a different capacity (corresponding to the capacity of the knapsack).

We implemented this algorithm in the code as an additional step to be run
before solving the kernel problem and the bucket sub-problems.
The testing however revealed that the algorithm couldn't find any valid bin packing
for any test instance.
After some consideration, we supposed that even if we were to swap FFD with another
heuristic algorithm, it would be unlikely that the bin packing found (if any) would impact the efficiency of the
Kernel Search by much.
Because of this reason, we decided to discard this technique.


\section{Single Knapsack Heuristic}
As an experiment, we tried to implement a heuristic algorithm that,
for each knapsack, solves a \textit{0-1 knapsack problem},
which finds the optimum items to be included in that knapsack.

Given \(n\) items with a weight \(w_{j}\) and a profit \(p_{j}\),
and a knapsack with capacity \(c\), the \textit{0-1 Knapsack Problem} finds
the subset of items that maximizes the overall profit while not overflowing the capacity of the knapsack.
\begin{equation}
    \max{z} = \sum_{j=1}^{n} p_{j} x_{j}
\end{equation}
\begin{equation}
    \sum_{j=1}^{n} w_{j} x_{j} \leq c
\end{equation}
\begin{equation}
    x_{j} \in \{0,1\} \qquad j=1,\dots,n
\end{equation}
It's worth noting that the \textit{0-1 Knapsack Problem} is equal to an MKP
where \(m=1\);

The detailed algorithm begins with defining \(J \coloneqq \text{items of the MKP}\)
and \(X \coloneqq \emptyset\) as the (initially empty) solution found by the heuristic.
Then, for each knapsack \(i = 1,\dots,m\) the following steps are repeated:
\begin{enumerate}
    \item Solve the \textit{0-1 Knapsack Problem} that uses \(i\) as the knapsack and the items in \(J\);
    \item Remove the items that were inserted into the knapsack from \(J\);
    \item Add the solution found (associated to the knapsack) to \(X\);
\end{enumerate}
Since every item can at most be inserted into one knapsack, and the capacity of each
knapsack is respected, \(X\) is guaranteed to be a feasible solution for the MKP\@.

After preliminarily testing the heuristic by itself (outside the Kernel Search framework)
we found that the results were excellent, and required very little computation time.

\subsection{Integration with the Kernel Search}
To integrate the heuristic with the Kernel Search, our idea was to use the heuristic
to find a high quality starting solution, which the Kernel Search tries to improve upon.

To achieve this, we removed the solving of the LP relaxation, and instead
used the solution found by the heuristic to build the kernel set and the buckets.
The parameters of the Kernel Search and the methods added until now
did not require any change to work correctly.

It is possible that the Kernel Search does not find a solution that improves
upon the one found by the heuristic: this is still acceptable, because the
heuristic finds integer and feasible solutions.



    % \chapter{Testing}\label{ch:testing}
In this chapter we will report the results of the tests executed
on the instances introduced in ~\ref{subsec:inst},
to verify the performance of the tuned Kernel Search algorithm.


\section{Default Bucket Construction Criterion}
The objective of the first round of testing was to find the most efficient
combination of kernel construction criterion and variable sorting criterion,
using the default bucket construction criterion (non-overlapping buckets).
Certain variable sorting criteria work really well with certain kernel construction criteria,
but give bad results with others.

We discovered that the two sorting criteria that
worked best were to sort by \textit{value, profit and weight, breaking ties using the reduced cost},
and the random sorting.
Random sorting however, due to its nature, produces inconsistent results.

The tables from ~\ref{tab:ker_pos_val_prof_wei_RC}
to ~\ref{tab:ker_tre_val_pro_wei_RC} contain the test results of all the
kernel construction criteria introduced in ~\ref{sec:kernel_criteria}
combined with the sorting by value, profit, weight and reduced cost.

The complete test results, including the ones with other sorting criteria, can be found at
\url{https://github.com/Golino98/KernelSearchGolinoCottiBeatrice/tree/main/log/DefaultBucket}.

In general, the performance of the criteria is comparable,
with the integer kernel being slightly more reliable.

\input{tables/kernelBase/KernelBuilderPositive/VariableSorterByValueProfitWeightAndRC}
\input{tables/kernelBase/KernelBuilderPercentage/VariableSorterByValueProfitWeightAndRC}
\input{tables/kernelBase/KernelBuilderIntValues/VariableSorterByValueProfitWeightAndRC}
\input{tables/kernelBase/KernelBuilderThreshold/VariableSorterByValueProfitWeightAndRC}


\section{Overlapping buckets}
In this section we report the results of the tests executed with the overlapping buckets.
Just like in the section before, we executed the tests for each possible combination
of kernel builder criterion and variable sorting criterion.

Tables from ~\ref{tab:ker_pos_val_prof_wei_RC_OVERL} to ~\ref{tab:ker_tre_val_pro_wei_RC_OVERL}
represent the results of the tests run with the various kernel builders,
and the variables sorted by \textit{value, profit and weight, breaking ties using the reduced cost}.
All the other tests results can be found at
\url{https://github.com/Golino98/KernelSearchGolinoCottiBeatrice/tree/main/log/BucketOverlap}.

The performance of the overlapping buckets on average slightly better (but still very much comparable)
to the non-overlapping ones.

\input{tables/bucketOverlap/KernelBuilderPositive/VariableSorterByValueProfitWeightAndRC}
\input{tables/bucketOverlap/KernelBuilderPercentage/VariableSorterByValueProfitWeightAndRC}
\input{tables/bucketOverlap/KernelBuilderIntValues/VariableSorterByValueProfitWeightAndRC}
\input{tables/bucketOverlap/KernelBuilderThreshold/VariableSorterByValueProfitWeightAndRC}


\section{Specific Improvements}
Using the results of the test above, we were able to determine
for each kernel construction criterion the most efficient
combination of variable sorting criterion and bucket
construction criterion (overlapping or non-overlapping buckets).

In the following sections, we report the performance
of each improvement explained in chapter ~\ref{ch:improvements},
using the best possible configuration for each kernel construction criterion.

To simplify the testing, each improvement was tested by itself: this means that a test suite was run for each
improvement, where all the other improvements are disabled.

Once again, not all test results are reported: the complete set can be found at \url{https://github
.com/Golino98/KernelSearchGolinoCottiBeatrice/tree/main/log}.

\subsection{Integer Kernel}
The best configuration for this kernel construction criterion is to use overlapping bucket,
and to sort variables by value, profit, weight and reduced cost.

\subsection{Percentage Kernel}
For this kernel construction criterion the best configuration uses overlapping buckets,
and sorts variables by reduced cost and value.
In general, we have noticed that this configuration behaves extremely well for small instances, but quite poorly on
bigger ones.

\subsection{Positive Kernel}
We have found three equally good configurations for this kernel construction criterion.
The one that gave slightly better results is the one that uses non-overlapping buckets and the random sorter.
It is worth noting however that, since variables are sorted randomly, the results may vary with each execution.

\subsection{Threshold Kernel}
The best configuration for the kernel construction that select variables based on a threshold is the one that uses
overlapping buckets and sort by value, profit, weight and reduced cost.

\input{tables/comparisonTables/Eject}
\input{tables/comparisonTables/RepCounter}
\input{tables/comparisonTables/ItemDom}
\input{tables/comparisonTables/Heuristic}

    % \chapter{Conclusions}
In this report we have described various improvements that we developed to
better tune the Kernel Search algorithm to the MKP\@.

Some of these improvements, such as Item Dominance and the Single Knapsack Heuristic
improve the performance and quality of the algorithm,
and thus can be seen as an improvement across the board.
Others, such as the Instance reduction and the Repetition Counter
usually worsen the results obtained but, depending on how they are configured,
they allow to find more solutions quickly and to diversificate the solution space.



    \backmatter
    %\printbibliography[heading=bibintoc] % Inserisci bibliografia nell'indice
\end{document}